\section{AIME I 1986}
        \begin{enumerate}
        	\item (\textit{AIME I 1986}) What is the sum of the solutions to the equation $\sqrt[4]{x} = \frac{12}{7 - \sqrt[4]{x}}$ ?


 

	\item (\textit{AIME I 1986}) Evaluate the product 
 \[\left(\sqrt{5}+\sqrt{6}+\sqrt{7}\right)\left(\sqrt{5}+\sqrt{6}-\sqrt{7}\right)\left(\sqrt{5}-\sqrt{6}+\sqrt{7}\right)\left(-\sqrt{5}+\sqrt{6}+\sqrt{7}\right).\]


 

	\item (\textit{AIME I 1986}) If $\tan x+\tan y=25$ and $\cot x + \cot y=30$, what is $\tan(x+y)$?


 

	\item (\textit{AIME I 1986}) Determine $3x_4+2x_5$ if $x_1$, $x_2$, $x_3$, $x_4$, and $x_5$ satisfy the system of equations below.


 $2x_1+x_2+x_3+x_4+x_5=6$$x_1+2x_2+x_3+x_4+x_5=12$$x_1+x_2+2x_3+x_4+x_5=24$$x_1+x_2+x_3+2x_4+x_5=48$$x_1+x_2+x_3+x_4+2x_5=96$

	\item (\textit{AIME I 1986}) What is that largest \href{https://artofproblemsolving.com/wiki/index.php?title=Positive\_integer}{positive integer} $n$ for which $n^3+100$ is \href{https://artofproblemsolving.com/wiki/index.php?title=Divisible}{divisible} by $n+10$?


 

	\item (\textit{AIME I 1986}) The pages of a book are numbered $1$ through $n$. When the page numbers of the book were added, one of the page numbers was mistakenly added twice, resulting in an incorrect sum of $1986$. What was the number of the page that was added twice?


 

	\item (\textit{AIME I 1986}) The increasing sequence $1,3,4,9,10,12,13\cdots$ consists of all those positive integers which are \href{https://artofproblemsolving.com/wiki/index.php?title=Powers&action=edit&redlink=1}{powers} of 3 or sums of distinct powers of 3. Find the $100^{\mbox{th}}$ term of this sequence.


 

	\item (\textit{AIME I 1986}) Let $S$ be the sum of the base $10$ logarithms of all the proper divisors of $1{,}000{,}000$. What is the integer nearest to $S$?


 

	\item (\textit{AIME I 1986}) In $\triangle ABC$, $AB= 425$, $BC=450$, and $AC=510$. An interior point $P$ is then drawn, and segments are drawn through $P$ parallel to the sides of the triangle. If these three segments are of an equal length $d$, find $d$.


 

	\item (\textit{AIME I 1986}) In a parlor game, the magician asks one of the participants to think of a three digit number $(abc)$ where $a$, $b$, and $c$ represent digits in base $10$ in the order indicated. The magician then asks this person to form the numbers $(acb)$, $(bca)$, $(bac)$, $(cab)$, and $(cba)$, to add these five numbers, and to reveal their sum, $N$. If told the value of $N$, the magician can identify the original number, $(abc)$. Play the role of the magician and determine the $(abc)$ if $N= 3194$.


 

	\item (\textit{AIME I 1986}) The polynomial $1-x+x^2-x^3+\cdots+x^{16}-x^{17}$ may be written in the form $a_0+a_1y+a_2y^2+\cdots +a_{16}y^{16}+a_{17}y^{17}$, where $y=x+1$ and the $a_i$'s are constants. Find the value of $a_2$.


 

	\item (\textit{AIME I 1986}) Let the sum of a set of numbers be the sum of its elements. Let $S$ be a set of positive integers, none greater than 15. Suppose no two disjoint subsets of $S$ have the same sum. What is the largest sum a set $S$ with these properties can have?


 

	\item (\textit{AIME I 1986}) In a sequence of coin tosses, one can keep a record of instances in which a tail is immediately followed by a head, a head is immediately followed by a head, and etc. We denote these by TH, HH, and etc. For example, in the sequence TTTHHTHTTTHHTTH of 15 coin tosses we observe that there are two HH, three HT, four TH, and five TT subsequences. How many different sequences of 15 coin tosses will contain exactly two HH, three HT, four TH, and five TT subsequences?


 

	\item (\textit{AIME I 1986}) The shortest distances between an interior \href{https://artofproblemsolving.com/wiki/index.php?title=Diagonal}{diagonal} of a rectangular \href{https://artofproblemsolving.com/wiki/index.php?title=Parallelepiped}{parallelepiped}, $P$, and the edges it does not meet are $2\sqrt{5}$, $\frac{30}{\sqrt{13}}$, and $\frac{15}{\sqrt{10}}$. Determine the \href{https://artofproblemsolving.com/wiki/index.php?title=Volume}{volume} of $P$.


 

	\item (\textit{AIME I 1986}) Let triangle $ABC$ be a right triangle in the $xy$-plane with a right angle at $C_{}$. Given that the length of the hypotenuse $AB$ is $60$, and that the medians through $A$ and $B$ lie along the lines $y=x+3$ and $y=2x+4$ respectively, find the area of triangle $ABC$.


 

\end{enumerate}